% TikZ 图：光电效应实验装置
% Day 27：近代物理
\begin{tikzpicture}[scale=0.9,>=Latex]
  \draw[thick,rounded corners] (-3,-1.4) rectangle (3,1.4);
  \node[above] at (0,1.4) {抽成真空的光电管};
  \draw[thick] (-1.8,-0.8) -- (-1.8,0.8);
  \node[below] at (-1.8,-0.82) {阴极 K};
  \draw[thick] (1.8,-0.8) -- (1.8,0.8);
  \node[right] at (1.8,0) {阳极 A};
  \foreach \y in {-0.45,0,0.45} {
    \draw[->,blue,thick] (-3.55,\y+0.35) -- (-1.95,\y);
  }
  \node[blue] at (-3.75,1.05) {入射光};
  \draw[->,red,thick] (-1.6,0.2) -- (1.45,0.45) node[midway,above] {光电子};
  \draw[->,red,thick] (-1.6,-0.25) -- (1.45,-0.15);
  \draw[thick] (-1.8,-1.4) -- (-1.8,-2.0) -- (-0.95,-2.0);
  \draw[thick] (-0.70,-2.0) circle (0.25);
  \node at (-0.70,-2.0) {$A$};
  \draw[thick] (-0.45,-2.0) -- (0.72,-2.0);
  \draw[thick] (0.72,-2.30) -- (0.72,-1.70);
  \draw[thick] (1.02,-2.16) -- (1.02,-1.84);
  \draw[thick] (1.02,-2.0) -- (1.8,-2.0) -- (1.8,-1.4);
  \node[font=\scriptsize] at (-1.35,-2.45) {微安表};
  \node[font=\scriptsize] at (1.45,-2.45) {可调电源（可反接）};
  \node[font=\small,align=center] at (0,-2.92)
    {调节外加电压，观察光电流是否存在以及何时为零};
\end{tikzpicture}
